% ============================================================================
%  PLANTILLA — GUÍA TÉCNICA EDUCATIVA PARA ASIGNATURAS DE CÓDIGO
%  Salida: PDF vía pdflatex (compilar dos veces para fijar referencias/índice)
%
%  Cómo usar esta plantilla:
%   1. NO edites el preámbulo salvo metadatos y el lenguaje de \lstset.
%   2. Reemplaza el contenido entre marcas << ... >> con el material real.
%   3. Cada bloque está comentado con su función pedagógica.
%   4. Compila:  pdflatex plantilla-guia.tex  (x2)
% ============================================================================

\documentclass[11pt,a4paper]{article}

% ---- Idioma y codificación --------------------------------------------------
\usepackage[utf8]{inputenc}
\usepackage[T1]{fontenc}
\usepackage[spanish,es-noquoting]{babel}

% ---- Layout, tipografía, color ---------------------------------------------
\usepackage[a4paper,margin=2.5cm,headheight=14pt]{geometry}
\usepackage{xcolor}
\usepackage{amsmath}
\usepackage{enumitem}
\usepackage{fancyhdr}
\usepackage{titlesec}
\usepackage[colorlinks=true,linkcolor=azulGuia,urlcolor=azulGuia,citecolor=azulGuia]{hyperref}

% ---- Paleta -----------------------------------------------------------------
\definecolor{azulGuia}{HTML}{1F4E79}
\definecolor{grisCodigo}{HTML}{F5F5F5}
\definecolor{verdeComentario}{HTML}{2E7D32}
\definecolor{moradoKeyword}{HTML}{6A1B9A}
\definecolor{rojoString}{HTML}{C62828}
\definecolor{ambar}{HTML}{B45309}

% ---- Cajas (objetivos, definición, ejercicio, etc.) ------------------------
\usepackage[most]{tcolorbox}

\newtcolorbox{objetivos}{colback=azulGuia!6,colframe=azulGuia,
  fonttitle=\bfseries,title=Objetivos de aprendizaje,arc=2mm,boxrule=0.6pt}
\newtcolorbox{definicion}[1]{colback=white,colframe=azulGuia,
  fonttitle=\bfseries,title=Definición: #1,arc=2mm,boxrule=0.6pt}
\newtcolorbox{nota}{colback=verdeComentario!8,colframe=verdeComentario,
  fonttitle=\bfseries,title=Nota,arc=2mm,boxrule=0.6pt}
\newtcolorbox{advertencia}{colback=ambar!10,colframe=ambar,
  fonttitle=\bfseries,title=Cuidado,arc=2mm,boxrule=0.6pt}
\newtcolorbox{ejemplo}[1]{colback=grisCodigo,colframe=azulGuia!70,
  fonttitle=\bfseries,title=Ejemplo resuelto: #1,arc=2mm,boxrule=0.6pt}

% Contador y caja para ejercicios numerados automáticamente
\newcounter{ejcontador}
\newtcolorbox{ejercicio}[1][]{colback=white,colframe=azulGuia!55,
  fonttitle=\bfseries,arc=2mm,boxrule=0.6pt,
  title={Ejercicio \stepcounter{ejcontador}\theejcontador\ifx\relax#1\relax\else\ (#1)\fi}}

% ---- Resaltado de código (listings) ----------------------------------------
% Cambia "language=Python" por C, C++, Java, SQL, etc. para otra asignatura.
\usepackage{listings}
\lstset{
  language=Python,
  backgroundcolor=\color{grisCodigo},
  basicstyle=\ttfamily\small,
  keywordstyle=\color{moradoKeyword}\bfseries,
  commentstyle=\color{verdeComentario}\itshape,
  stringstyle=\color{rojoString},
  numbers=left,
  numberstyle=\tiny\color{gray},
  numbersep=8pt,
  showstringspaces=false,
  breaklines=true,
  breakatwhitespace=true,
  frame=single,
  framesep=4pt,
  rulecolor=\color{gray!40},
  xleftmargin=14pt,
  tabsize=4,
  literate={á}{{\'a}}1 {é}{{\'e}}1 {í}{{\'i}}1 {ó}{{\'o}}1 {ú}{{\'u}}1
           {ñ}{{\~n}}1 {Á}{{\'A}}1 {É}{{\'E}}1 {Í}{{\'I}}1 {Ó}{{\'O}}1
           {Ú}{{\'U}}1 {Ñ}{{\~N}}1
}

% ---- Encabezado y pie -------------------------------------------------------
\pagestyle{fancy}
\fancyhf{}
\lhead{\small\color{azulGuia}<<NOMBRE DE LA ASIGNATURA>>}
\rhead{\small\color{azulGuia}<<TEMA>>}
\cfoot{\small Página \thepage}
\renewcommand{\headrulewidth}{0.4pt}

% ---- Estilo de títulos ------------------------------------------------------
\titleformat{\section}{\Large\bfseries\color{azulGuia}}{\thesection}{1em}{}
\titleformat{\subsection}{\large\bfseries\color{azulGuia!85}}{\thesubsection}{1em}{}

% ============================================================================
%  METADATOS DE PORTADA
% ============================================================================
\title{\vspace{-1cm}\color{azulGuia}\textbf{<<TÍTULO DE LA GUÍA>>}\\[2mm]
  \large\normalcolor <<NOMBRE DE LA ASIGNATURA>> --- Unidad <<N>>}
\author{<<Docente / Institución>>}
\date{<<Semestre / Fecha>>}

\begin{document}
\maketitle
\thispagestyle{fancy}

% ============================================================================
%  1. OBJETIVOS  — qué sabrá hacer el estudiante al terminar (verbos medibles)
% ============================================================================
\begin{objetivos}
Al finalizar esta guía, el estudiante será capaz de:
\begin{itemize}[leftmargin=1.4em,nosep]
  \item Explicar qué es la recursividad y distinguir caso base de caso recursivo.
  \item Trazar la pila de llamadas de una función recursiva paso a paso.
  \item Implementar soluciones recursivas correctas en Python.
  \item Reconocer cuándo la recursión es preferible a la iteración y viceversa.
\end{itemize}
\end{objetivos}

% ============================================================================
%  PRERREQUISITOS  — qué debe saber el estudiante de antemano
% ============================================================================
\section*{Conocimientos previos}
Antes de comenzar, debes manejar: definición de funciones en Python,
parámetros y valores de retorno, condicionales (\texttt{if/else}) y el concepto
de variable local.

% ============================================================================
%  2. DESARROLLO CONCEPTUAL  — teoría progresiva, de lo simple a lo complejo
% ============================================================================
\section{¿Qué es la recursividad?}

Una función \textbf{recursiva} es aquella que se llama a sí misma para resolver
una versión más pequeña del mismo problema. La idea descansa en dos piezas que
deben coexistir siempre.

\begin{definicion}{recursividad}
Técnica en la que la solución de un problema se expresa en términos de
soluciones de instancias más pequeñas del mismo problema, hasta alcanzar un
\emph{caso base} que se resuelve directamente.
\end{definicion}

\subsection{Las dos piezas obligatorias}

\begin{enumerate}[leftmargin=1.6em]
  \item \textbf{Caso base:} la condición de parada. Sin él, la función se llama
        infinitamente y agota la pila.
  \item \textbf{Caso recursivo:} la llamada a sí misma con una entrada
        estrictamente más cercana al caso base.
\end{enumerate}

\begin{advertencia}
Si el caso recursivo no \emph{avanza} hacia el caso base (por ejemplo, llamar
con el mismo argumento), obtendrás un \texttt{RecursionError} en Python.
\end{advertencia}

\subsection{Cómo razona la máquina: la pila de llamadas}

Cada llamada pendiente se apila; cuando un caso base retorna, la pila se
desapila resolviendo cada nivel. Comprender este apilamiento es la clave para
depurar recursión.

% ============================================================================
%  3. EJEMPLOS RESUELTOS  — código real + explicación línea a línea
% ============================================================================
\section{Ejemplos resueltos}

\begin{ejemplo}{factorial de un número}
El factorial $n! = n \times (n-1) \times \cdots \times 1$, con $0! = 1$.
\end{ejemplo}

\begin{lstlisting}
def factorial(n):
    # Caso base: el factorial de 0 es 1
    if n == 0:
        return 1
    # Caso recursivo: n * factorial de (n-1)
    return n * factorial(n - 1)

print(factorial(4))  # -> 24
\end{lstlisting}

\noindent\textbf{Traza de \texttt{factorial(4)}:}
\[
4 \cdot \underbrace{3 \cdot \underbrace{2 \cdot \underbrace{1 \cdot \underbrace{1}_{n=0}}_{n=1}}_{n=2}}_{n=3} = 24
\]

\begin{nota}
Observa que cada llamada reduce \texttt{n} en uno: así se garantiza que el
caso base ($n=0$) siempre se alcanza para entradas no negativas.
\end{nota}

% ============================================================================
%  4. EJERCICIOS  — graduados de menor a mayor dificultad
% ============================================================================
\section{Ejercicios propuestos}

\begin{ejercicio}[básico]
Escribe una función recursiva \texttt{suma\_hasta(n)} que devuelva
$1 + 2 + \cdots + n$. Identifica explícitamente el caso base.
\end{ejercicio}

\begin{ejercicio}[intermedio]
Implementa \texttt{potencia(base, exp)} de forma recursiva, sin usar el
operador \texttt{**}. Asume \texttt{exp} entero no negativo.
\end{ejercicio}

\begin{ejercicio}[desafío]
La sucesión de Fibonacci se define $F(0)=0$, $F(1)=1$,
$F(n)=F(n-1)+F(n-2)$. Impleméntala recursivamente y luego responde:
¿por qué esta versión es ineficiente para $n$ grande? Justifica en términos
del número de llamadas.
\end{ejercicio}

% ============================================================================
%  5. SOLUCIONES  — al final para no revelarlas antes de tiempo
% ============================================================================
\newpage
\section*{Soluciones}

\textbf{Ejercicio 1.}
\begin{lstlisting}
def suma_hasta(n):
    if n == 0:        # caso base
        return 0
    return n + suma_hasta(n - 1)
\end{lstlisting}

\textbf{Ejercicio 2.}
\begin{lstlisting}
def potencia(base, exp):
    if exp == 0:      # caso base: cualquier base^0 = 1
        return 1
    return base * potencia(base, exp - 1)
\end{lstlisting}

\textbf{Ejercicio 3.} La versión directa recalcula los mismos valores muchas
veces: \texttt{F(n)} genera dos llamadas, cada una otras dos, etc., dando un
crecimiento exponencial $O(2^n)$ de llamadas. Se resuelve con memoización o
con una versión iterativa $O(n)$.

% ============================================================================
%  REFERENCIAS
% ============================================================================
\section*{Referencias}
\begin{itemize}[leftmargin=1.4em,nosep]
  \item <<Autor, Título, Editorial, Año / URL con fecha de consulta>>
\end{itemize}

\end{document}
